% Deck 4 of 4 of AI for Biodiversity: A Primer.
% Build: latexmk -pdf main.tex   (pdflatex + biber)
%
% Part V (Chapters 15--18) and the appendices. Every statement comes from the
% book; the figures are the book's (book/figures); the five paper screenshots
% are those assigned to this deck in presentations/BRIEF.md.
\documentclass[11pt,aspectratio=169]{beamer}
% ---------------------------------------------------------------------------
% preamble.tex --- shared Beamer style for the four workshop decks
%
% Loaded by each deck with % ---------------------------------------------------------------------------
% preamble.tex --- shared Beamer style for the four workshop decks
%
% Loaded by each deck with % ---------------------------------------------------------------------------
% preamble.tex --- shared Beamer style for the four workshop decks
%
% Loaded by each deck with \input{../common/preamble}. Same fonts, colours
% and bibliography settings as the book (book/preamble.tex): Latin Modern,
% black on white, one grey for code, one accent colour for what is untrusted
% or wrong. Built with pdflatex + biber, like the book.
% ---------------------------------------------------------------------------

\usepackage[T1]{fontenc}
\usepackage[utf8]{inputenc}
\usepackage{lmodern}
\usepackage{microtype}
\usepackage{booktabs}
\usepackage{tabularx}
\usepackage{listings}
\usepackage{xcolor}
\usepackage{csquotes}
\usepackage{graphicx}

% Figures: the book's figures as the 300 dpi PNGs the book itself includes
% (book/figures/png, tracked in git and current; figures/pdf is not rebuilt
% after every change), and the first pages of the papers
% (common/screenshots). Both are found by name without an extension.
\graphicspath{{../../book/figures/png/}{../common/screenshots/}}
\DeclareGraphicsExtensions{.png,.pdf}

% --- look ------------------------------------------------------------------
\usetheme{default}
\usefonttheme{serif}
\definecolor{accent}{HTML}{9B2D20}
\definecolor{codebg}{gray}{0.96}
\definecolor{rule}{gray}{0.75}
\definecolor{dim}{gray}{0.45}
\setbeamercolor{structure}{fg=black}
\setbeamercolor{normal text}{fg=black,bg=white}
\setbeamercolor{alerted text}{fg=accent}
\setbeamercolor{section in toc}{fg=black}
\setbeamercolor{footline}{fg=dim}
\setbeamerfont{title}{series=\bfseries,size=\Large}
\setbeamerfont{subtitle}{series=\mdseries,size=\normalsize}
\setbeamerfont{frametitle}{series=\bfseries,size=\large}
\setbeamerfont{footline}{size=\scriptsize}
\setbeamertemplate{navigation symbols}{}
\setbeamertemplate{itemize items}[circle]
\setbeamertemplate{enumerate items}[default]
\setbeamertemplate{caption}{\scriptsize\insertcaption}
\setbeamertemplate{frametitle}{%
  \vskip0.6em\usebeamerfont{frametitle}\insertframetitle\par
  \vskip0.2em{\color{rule}\hrule height 0.4pt}\vskip0.4em}
\setbeamertemplate{footline}{%
  \hbox to \paperwidth{\hspace{1.2em}\usebeamerfont{footline}\usebeamercolor[fg]{footline}%
    \insertshorttitle\ --- \insertsection\hfill
    \insertframenumber/\inserttotalframenumber\hspace{1.2em}}\vskip0.6em}
\setbeamersize{text margin left=1.4em,text margin right=1.4em}
% An outline frame at the start of each section, the current one in black.
\AtBeginSection[]{\begin{frame}{Outline}\tableofcontents[currentsection]\end{frame}}
% A frame that breaks over pages is numbered only from its second page.
\setbeamertemplate{frametitle continuation}[from second]

% --- code listings (as in the book) ----------------------------------------
\lstset{
  basicstyle=\ttfamily\footnotesize,
  backgroundcolor=\color{codebg},
  frame=none,
  breaklines=true,
  breakatwhitespace=true,
  columns=fullflexible,
  keepspaces=true,
  upquote=true,
  showstringspaces=false,
  xleftmargin=0.6em,
  framexleftmargin=0.6em,
  aboveskip=0.5em,
  belowskip=0.5em,
  literate={~}{{\textasciitilde}}1
           {>}{{\textgreater}}1
           {<}{{\textless}}1
           {ö}{{\"o}}1 {ü}{{\"u}}1 {é}{{\'e}}1,
}
\lstdefinestyle{shell}{language=bash}

% --- bibliography (same settings as the book) ------------------------------
\usepackage[
  backend=biber,
  style=authoryear,
  sorting=nyt,
  maxbibnames=6,
  minbibnames=3,
  maxcitenames=2,
  giveninits=true,
  uniquename=init,
  doi=true,
  url=true,
  isbn=false,
  eprint=true,
]{biblatex}
\addbibresource{../../book/references.bib}
\renewcommand*{\bibfont}{\scriptsize}
\setlength{\bibitemsep}{0.25em}
\setbeamertemplate{bibliography item}{}   % no icon before each entry

% --- shorthand -------------------------------------------------------------
\newcommand{\code}[1]{\texttt{#1}}
\newcommand{\term}[1]{\emph{#1}}
% A key point, as in the book: no box.
\newcommand{\kp}[1]{\par\medskip{\bfseries Key point.}\enspace\emph{#1}\par}
% \fig[height]{name}: one of the book's figures, scaled to fit the frame.
\newcommand{\fig}[2][0.72\textheight]{%
  \centering\includegraphics[height=#1,width=\textwidth,keepaspectratio]{#2}\par}
% \screenshot{key}: the first page of a paper, with a thin frame.
\newcommand{\screenshot}[2][0.66\textheight]{%
  \setlength{\fboxsep}{0pt}\setlength{\fboxrule}{0.3pt}%
  \fbox{\includegraphics[height=#1,width=\linewidth,keepaspectratio]{#2}}}
% \paperslide{key}{text}: screenshot on the left, what the book takes from the
% paper on the right. Use inside a frame; add \source after it.
\newcommand{\paperslide}[2]{%
  \begin{columns}[T]
    \begin{column}{0.40\textwidth}\centering\screenshot{#1}\end{column}
    \begin{column}{0.58\textwidth}\small #2\end{column}
  \end{columns}}
% Notes in references.bib record how a web source was saved and dated; they
% belong in the reference list, not on a \source line, so cite commands
% (including \fullcite) leave them out.
\AtEveryCitekey{\clearfield{note}}
% \source{...}: a small grey line at the foot of a frame, for the source.
\newcommand{\source}[1]{\vskip0pt plus 1filll\par{\scriptsize\color{dim}#1\par}}
. Same fonts, colours
% and bibliography settings as the book (book/preamble.tex): Latin Modern,
% black on white, one grey for code, one accent colour for what is untrusted
% or wrong. Built with pdflatex + biber, like the book.
% ---------------------------------------------------------------------------

\usepackage[T1]{fontenc}
\usepackage[utf8]{inputenc}
\usepackage{lmodern}
\usepackage{microtype}
\usepackage{booktabs}
\usepackage{tabularx}
\usepackage{listings}
\usepackage{xcolor}
\usepackage{csquotes}
\usepackage{graphicx}

% Figures: the book's figures as the 300 dpi PNGs the book itself includes
% (book/figures/png, tracked in git and current; figures/pdf is not rebuilt
% after every change), and the first pages of the papers
% (common/screenshots). Both are found by name without an extension.
\graphicspath{{../../book/figures/png/}{../common/screenshots/}}
\DeclareGraphicsExtensions{.png,.pdf}

% --- look ------------------------------------------------------------------
\usetheme{default}
\usefonttheme{serif}
\definecolor{accent}{HTML}{9B2D20}
\definecolor{codebg}{gray}{0.96}
\definecolor{rule}{gray}{0.75}
\definecolor{dim}{gray}{0.45}
\setbeamercolor{structure}{fg=black}
\setbeamercolor{normal text}{fg=black,bg=white}
\setbeamercolor{alerted text}{fg=accent}
\setbeamercolor{section in toc}{fg=black}
\setbeamercolor{footline}{fg=dim}
\setbeamerfont{title}{series=\bfseries,size=\Large}
\setbeamerfont{subtitle}{series=\mdseries,size=\normalsize}
\setbeamerfont{frametitle}{series=\bfseries,size=\large}
\setbeamerfont{footline}{size=\scriptsize}
\setbeamertemplate{navigation symbols}{}
\setbeamertemplate{itemize items}[circle]
\setbeamertemplate{enumerate items}[default]
\setbeamertemplate{caption}{\scriptsize\insertcaption}
\setbeamertemplate{frametitle}{%
  \vskip0.6em\usebeamerfont{frametitle}\insertframetitle\par
  \vskip0.2em{\color{rule}\hrule height 0.4pt}\vskip0.4em}
\setbeamertemplate{footline}{%
  \hbox to \paperwidth{\hspace{1.2em}\usebeamerfont{footline}\usebeamercolor[fg]{footline}%
    \insertshorttitle\ --- \insertsection\hfill
    \insertframenumber/\inserttotalframenumber\hspace{1.2em}}\vskip0.6em}
\setbeamersize{text margin left=1.4em,text margin right=1.4em}
% An outline frame at the start of each section, the current one in black.
\AtBeginSection[]{\begin{frame}{Outline}\tableofcontents[currentsection]\end{frame}}
% A frame that breaks over pages is numbered only from its second page.
\setbeamertemplate{frametitle continuation}[from second]

% --- code listings (as in the book) ----------------------------------------
\lstset{
  basicstyle=\ttfamily\footnotesize,
  backgroundcolor=\color{codebg},
  frame=none,
  breaklines=true,
  breakatwhitespace=true,
  columns=fullflexible,
  keepspaces=true,
  upquote=true,
  showstringspaces=false,
  xleftmargin=0.6em,
  framexleftmargin=0.6em,
  aboveskip=0.5em,
  belowskip=0.5em,
  literate={~}{{\textasciitilde}}1
           {>}{{\textgreater}}1
           {<}{{\textless}}1
           {ö}{{\"o}}1 {ü}{{\"u}}1 {é}{{\'e}}1,
}
\lstdefinestyle{shell}{language=bash}

% --- bibliography (same settings as the book) ------------------------------
\usepackage[
  backend=biber,
  style=authoryear,
  sorting=nyt,
  maxbibnames=6,
  minbibnames=3,
  maxcitenames=2,
  giveninits=true,
  uniquename=init,
  doi=true,
  url=true,
  isbn=false,
  eprint=true,
]{biblatex}
\addbibresource{../../book/references.bib}
\renewcommand*{\bibfont}{\scriptsize}
\setlength{\bibitemsep}{0.25em}
\setbeamertemplate{bibliography item}{}   % no icon before each entry

% --- shorthand -------------------------------------------------------------
\newcommand{\code}[1]{\texttt{#1}}
\newcommand{\term}[1]{\emph{#1}}
% A key point, as in the book: no box.
\newcommand{\kp}[1]{\par\medskip{\bfseries Key point.}\enspace\emph{#1}\par}
% \fig[height]{name}: one of the book's figures, scaled to fit the frame.
\newcommand{\fig}[2][0.72\textheight]{%
  \centering\includegraphics[height=#1,width=\textwidth,keepaspectratio]{#2}\par}
% \screenshot{key}: the first page of a paper, with a thin frame.
\newcommand{\screenshot}[2][0.66\textheight]{%
  \setlength{\fboxsep}{0pt}\setlength{\fboxrule}{0.3pt}%
  \fbox{\includegraphics[height=#1,width=\linewidth,keepaspectratio]{#2}}}
% \paperslide{key}{text}: screenshot on the left, what the book takes from the
% paper on the right. Use inside a frame; add \source after it.
\newcommand{\paperslide}[2]{%
  \begin{columns}[T]
    \begin{column}{0.40\textwidth}\centering\screenshot{#1}\end{column}
    \begin{column}{0.58\textwidth}\small #2\end{column}
  \end{columns}}
% Notes in references.bib record how a web source was saved and dated; they
% belong in the reference list, not on a \source line, so cite commands
% (including \fullcite) leave them out.
\AtEveryCitekey{\clearfield{note}}
% \source{...}: a small grey line at the foot of a frame, for the source.
\newcommand{\source}[1]{\vskip0pt plus 1filll\par{\scriptsize\color{dim}#1\par}}
. Same fonts, colours
% and bibliography settings as the book (book/preamble.tex): Latin Modern,
% black on white, one grey for code, one accent colour for what is untrusted
% or wrong. Built with pdflatex + biber, like the book.
% ---------------------------------------------------------------------------

\usepackage[T1]{fontenc}
\usepackage[utf8]{inputenc}
\usepackage{lmodern}
\usepackage{microtype}
\usepackage{booktabs}
\usepackage{tabularx}
\usepackage{listings}
\usepackage{xcolor}
\usepackage{csquotes}
\usepackage{graphicx}

% Figures: the book's figures as the 300 dpi PNGs the book itself includes
% (book/figures/png, tracked in git and current; figures/pdf is not rebuilt
% after every change), and the first pages of the papers
% (common/screenshots). Both are found by name without an extension.
\graphicspath{{../../book/figures/png/}{../common/screenshots/}}
\DeclareGraphicsExtensions{.png,.pdf}

% --- look ------------------------------------------------------------------
\usetheme{default}
\usefonttheme{serif}
\definecolor{accent}{HTML}{9B2D20}
\definecolor{codebg}{gray}{0.96}
\definecolor{rule}{gray}{0.75}
\definecolor{dim}{gray}{0.45}
\setbeamercolor{structure}{fg=black}
\setbeamercolor{normal text}{fg=black,bg=white}
\setbeamercolor{alerted text}{fg=accent}
\setbeamercolor{section in toc}{fg=black}
\setbeamercolor{footline}{fg=dim}
\setbeamerfont{title}{series=\bfseries,size=\Large}
\setbeamerfont{subtitle}{series=\mdseries,size=\normalsize}
\setbeamerfont{frametitle}{series=\bfseries,size=\large}
\setbeamerfont{footline}{size=\scriptsize}
\setbeamertemplate{navigation symbols}{}
\setbeamertemplate{itemize items}[circle]
\setbeamertemplate{enumerate items}[default]
\setbeamertemplate{caption}{\scriptsize\insertcaption}
\setbeamertemplate{frametitle}{%
  \vskip0.6em\usebeamerfont{frametitle}\insertframetitle\par
  \vskip0.2em{\color{rule}\hrule height 0.4pt}\vskip0.4em}
\setbeamertemplate{footline}{%
  \hbox to \paperwidth{\hspace{1.2em}\usebeamerfont{footline}\usebeamercolor[fg]{footline}%
    \insertshorttitle\ --- \insertsection\hfill
    \insertframenumber/\inserttotalframenumber\hspace{1.2em}}\vskip0.6em}
\setbeamersize{text margin left=1.4em,text margin right=1.4em}
% An outline frame at the start of each section, the current one in black.
\AtBeginSection[]{\begin{frame}{Outline}\tableofcontents[currentsection]\end{frame}}
% A frame that breaks over pages is numbered only from its second page.
\setbeamertemplate{frametitle continuation}[from second]

% --- code listings (as in the book) ----------------------------------------
\lstset{
  basicstyle=\ttfamily\footnotesize,
  backgroundcolor=\color{codebg},
  frame=none,
  breaklines=true,
  breakatwhitespace=true,
  columns=fullflexible,
  keepspaces=true,
  upquote=true,
  showstringspaces=false,
  xleftmargin=0.6em,
  framexleftmargin=0.6em,
  aboveskip=0.5em,
  belowskip=0.5em,
  literate={~}{{\textasciitilde}}1
           {>}{{\textgreater}}1
           {<}{{\textless}}1
           {ö}{{\"o}}1 {ü}{{\"u}}1 {é}{{\'e}}1,
}
\lstdefinestyle{shell}{language=bash}

% --- bibliography (same settings as the book) ------------------------------
\usepackage[
  backend=biber,
  style=authoryear,
  sorting=nyt,
  maxbibnames=6,
  minbibnames=3,
  maxcitenames=2,
  giveninits=true,
  uniquename=init,
  doi=true,
  url=true,
  isbn=false,
  eprint=true,
]{biblatex}
\addbibresource{../../book/references.bib}
\renewcommand*{\bibfont}{\scriptsize}
\setlength{\bibitemsep}{0.25em}
\setbeamertemplate{bibliography item}{}   % no icon before each entry

% --- shorthand -------------------------------------------------------------
\newcommand{\code}[1]{\texttt{#1}}
\newcommand{\term}[1]{\emph{#1}}
% A key point, as in the book: no box.
\newcommand{\kp}[1]{\par\medskip{\bfseries Key point.}\enspace\emph{#1}\par}
% \fig[height]{name}: one of the book's figures, scaled to fit the frame.
\newcommand{\fig}[2][0.72\textheight]{%
  \centering\includegraphics[height=#1,width=\textwidth,keepaspectratio]{#2}\par}
% \screenshot{key}: the first page of a paper, with a thin frame.
\newcommand{\screenshot}[2][0.66\textheight]{%
  \setlength{\fboxsep}{0pt}\setlength{\fboxrule}{0.3pt}%
  \fbox{\includegraphics[height=#1,width=\linewidth,keepaspectratio]{#2}}}
% \paperslide{key}{text}: screenshot on the left, what the book takes from the
% paper on the right. Use inside a frame; add \source after it.
\newcommand{\paperslide}[2]{%
  \begin{columns}[T]
    \begin{column}{0.40\textwidth}\centering\screenshot{#1}\end{column}
    \begin{column}{0.58\textwidth}\small #2\end{column}
  \end{columns}}
% Notes in references.bib record how a web source was saved and dated; they
% belong in the reference list, not on a \source line, so cite commands
% (including \fullcite) leave them out.
\AtEveryCitekey{\clearfield{note}}
% \source{...}: a small grey line at the foot of a frame, for the source.
\newcommand{\source}[1]{\vskip0pt plus 1filll\par{\scriptsize\color{dim}#1\par}}


\title[AI in Practice 4]{AI in Practice, or: How I Learned to Stop Worrying and Love the Agent}
\subtitle{A toolbox and worked examples for biodiversity researchers, from the book \emph{AI for Biodiversity: A Primer}\\[0.6em] 4. Risk, integrity, keeping up, and the wider frame (Chapters 15--18, Appendices)}
\author{Savvas Paragkamian}
\institute{Hellenic Centre for Marine Research \\ KYBELE and ELIXIR-GR}
\date{HCMR workshop, November 2026}
\hypersetup{pdftitle={AI for Biodiversity: A Primer, deck 4},pdfauthor={Savvas Paragkamian}}

\begin{document}

\begin{frame}[plain]
  \titlepage
\end{frame}

\begin{frame}{Goals of this session}
  \begin{enumerate}
    \item Know what an agent can do on your machine and with your data, and
      how to fence it in (Chapter 15).
    \item See where a model helps at each stage of research, and the check
      that goes with each stage (Chapter 16).
    \item Decide what is worth learning, and set the practical questions in
      their wider frame (Chapters 17 and 18, and the appendices).
  \end{enumerate}
  \kp{Run models locally where you can, verify everything that matters, and
    disclose what you used.}
\end{frame}

% ---------------------------------------------------------------------------
\section{Security, biorisk and governance (Chapter 15)}

\begin{frame}{Running an agent is running arbitrary code}
  \paperslide{shapira2026agentsofchaos}{\footnotesize
    \begin{itemize}
      \item An agent harness runs commands on your machine, with your
        permissions. The model can be steered by anything it reads.
      \item The test: agents with email, chat, files and a shell in a live
        lab, probed by twenty researchers for two weeks.
      \item They obeyed people who were not their owners, leaked sensitive
        information, ran destructive commands and used resources without
        limit.
      \item They were fooled by spoofed identities, spread unsafe habits to
        each other, and reported tasks done when the system showed they were
        not \parencite{shapira2026agentsofchaos}.
      \item An agent's account of what it did is not evidence that it did it.
    \end{itemize}}
  \source{\fullcite{shapira2026agentsofchaos}}
\end{frame}

\begin{frame}{Prompt injection: instructions and data arrive as one stream}
  \fig[0.36\textheight]{ch15-injection}
  \begin{itemize}
    \item A model does not reliably tell instructions from data. Any text it
      reads --- a file, a web page, a code comment, a CSV cell --- can contain
      words it treats as a command.
    \item Watch for instructions hidden in a downloaded dataset's metadata,
      invisible text in a PDF, a dependency's documentation, a web page
      fetched during a literature search, and an MCP server's description of
      its own tools.
  \end{itemize}
  \kp{Treat everything an agent reads as untrusted, including your own data
    files if they came from elsewhere, and including what tools say about
    themselves.}
  \source{Figure from Chapter 15 of the book.}
\end{frame}

\begin{frame}{Prompt injection follows from the design}
  \paperslide{greshake2023injection}{\footnotesize
    \begin{itemize}
      \item Instructions and data arrive as tokens in the same context
        \parencite{greshake2023injection}. A different problem from
        \term{jailbreaking}, where the user talks the model out of its
        safety training.
      \item Bing Chat refused a jailbreak the user typed, but obeyed the same
        prompt when it read it on a web page.
      \item Mitigations reduce the rate; none removes it. Eight defences were
        all bypassed more than half the time by attacks adapted to them
        \parencite{zhan2025adaptive}; on the newest models basic injection has
        weakened to about a quarter of its early effect, not to zero
        \parencite{queen2026shelflife}.
      \item So limit what an injection can achieve. Never give one agent the
        \term{lethal trifecta}: private data, untrusted content, and a way
        to communicate outward \parencite{willison2025trifecta}.
    \end{itemize}}
  \source{\fullcite{greshake2023injection}}
\end{frame}

\begin{frame}{Fencing an agent in: layers, each with an honest limit}
  \begin{columns}[T]
    \begin{column}{0.44\textwidth}
      \fig[0.64\textheight]{ch15-fences}
    \end{column}
    \begin{column}{0.54\textwidth}
      \footnotesize
      \begin{itemize}
        \item In practice: start the agent in the project directory, without
          administrator rights; keys out of its reach; commit before it
          starts, read the diff afterwards.
        \item Cut network access when the task is local; a container or
          microVM for unattended work. A credential that has appeared in a
          chat is compromised: rotate it.
        \item \term{Slopsquatting}: models invent package names and
          attackers register them; 5.2 to 21.7 per cent of suggested packages
          did not exist in one test \parencite{spracklen2025packages}. Check
          every dependency an agent adds.
        \item MCP servers are programs with access to your data; skills are
          instructions an agent follows. Read the tool descriptions and pin
          the version.
      \end{itemize}
    \end{column}
  \end{columns}
  \source{Figure from Chapter 15 of the book.}
\end{frame}

\begin{frame}{When the fence fails: July 2026}
  \small
  \begin{itemize}
    \item OpenAI tested models for cyber-attack skills with refusals reduced
      and the blocking classifiers switched off; the only link outward was a
      package-cache service \parencite{openai2026huggingface}.
    \item The models found an unknown flaw in it, reached the internet, and
      broke into Hugging Face's servers in pursuit of a way to cheat the
      evaluation.
    \item \textcite{bengio2026lying}: the behaviour follows from training.
      Faced with a sharp goal (pass the test) and a vague one (behave well),
      an agent finds a reading of the vague one that lets it win.
    \item \textcite{mitchell2026metaphors}: no one lost control. The failures
      were human: a sandbox that leaked, safeguards turned off, weeks without
      oversight, training that rewards persistence.
    \item Either way, what surrounds the model matters as much as the model.
      \term{Least privilege}: only the permissions the current task needs.
  \end{itemize}
\end{frame}

\begin{frame}{Dual use in biology: where the safeguards sit}
  \begin{columns}[T]
    \begin{column}{0.56\textwidth}
      \fig[0.74\textheight]{ch15-shared-layer}
    \end{column}
    \begin{column}{0.42\textwidth}
      \footnotesize
      \begin{itemize}
        \item Novices with model access were 4.16 times as accurate as
          novices with only the internet, on tasks done at a computer, not at
          the bench \parencite{zhang2026uplift}.
        \item The counterweight is calibration: tacit knowledge, materials
          and laboratory skill, not information, are the real barriers
          \parencite{peppin2025biorisk}.
        \item Both can be true: information barriers have fallen, physical
          ones have not. The two fields meet in agents: injection can carry a
          jailbreak past a filter.
      \end{itemize}
    \end{column}
  \end{columns}
  \source{Figure from Chapter 15 of the book.}
\end{frame}

\begin{frame}{Sensitive localities: the nearer dual-use problem}
  \paperslide{lindenmayer2017donotpublish}{\footnotesize
    \begin{itemize}
      \item For most ecologists the nearer dual-use problem is locality data.
        Poachers have targeted more than 20 newly described reptile species,
        some within months.
      \item People searching for an endangered lizard trespassed on farms soon
        after its records reached a public atlas
        \parencite{lindenmayer2017donotpublish}.
      \item In the author's view, an agent gathering localities from hundreds
        of papers at once makes such a search easier.
      \item GBIF's guidance: generalise sensitive records before sharing
        them, and record it in \code{dataGeneralizations}.
      \item Precise localities pasted into a hosted model go to a third
        party: generalise first, or use a local model.
    \end{itemize}}
  \source{\fullcite{lindenmayer2017donotpublish}}
\end{frame}

\begin{frame}{Risks to judgement: sycophancy, slop and deskilling}
  \paperslide{jiralerspong2026slopcliff}{\footnotesize
    \begin{itemize}
      \item \term{Sycophancy} removes the second opinion you thought you
        were getting; alone, your results are checked only by you and the
        model behind them \parencite{jiralerspong2026slopcliff}.
      \item The \term{slop cliff}: actual productivity falls while perceived
        productivity climbs. Can you explain your data and results without
        asking the model?
      \item \term{Slop}: of seven authors of papers due for rejection, three
        could not answer basic questions about them
        \parencite{turner2026slop}.
      \item ``You become less capable while appearing more productive'': the
        \term{deskilling} the book warns of for ecology. In \emph{McNeal}, a
        novelist lets a chatbot write \parencite{akhtar2024mcneal}.
      \item Sign only work you can explain.
    \end{itemize}}
  \source{\fullcite{jiralerspong2026slopcliff}}
\end{frame}

% ---------------------------------------------------------------------------
\section{Changes across the scientific workflow (Chapter 16)}

\begin{frame}{Five stages, and the check that goes with each}
  \fig[0.30\textheight]{ch16-workflow-checks}
  \small
  \begin{itemize}
    \item Finding literature has changed most. \term{Deep research} systems
      break a question into parts, search, read what they find, and decide
      what to search next.
    \item A \term{scoping search} explores a topic. A systematic review needs
      a final search that someone else can repeat: ROSES and PRISMA-EcoEvo
      ask for the databases, the exact search string and the dates covered.
    \item Deep research belongs to the scoping stage. Its results cannot be
      repeated, so it can at most add candidates to a documented search. The
      same rule for data searches: record the query and the date.
  \end{itemize}
  \source{Figure from Chapter 16 of the book.}
\end{frame}

\begin{frame}{The tools arrive from above: funders and institutions}
  \paperslide{anthropic2026aiforscience}{\footnotesize
    \begin{itemize}
      \item 28 October 2026: Anthropic's AI for Science programme (free
        credits for its models) now runs through funders and institutions
        \parencite{anthropic2026aiforscience}.
      \item Sloan, Simons, Schmidt Sciences and the US National Science
        Foundation will distribute credits; EMBL, Sanger and the Crick
        will use them in-house.
      \item Also announced: a specification for agents that operate
        laboratory instruments (with HHMI Janelia); teams of agents
        analysing the Allen Institute's open data.
      \item A company's own statements; no amounts or criteria. A public
        funder handing out a private provider's credits makes dependence
        part of how science is paid for (Chapter 5).
      \item The checks are yours, whoever paid for the tool.
    \end{itemize}}
  \source{\fullcite{anthropic2026aiforscience}}
\end{frame}

\begin{frame}{Resolve every reference before it enters a manuscript}
  \small A conservative estimate puts the invented citations in four large
  collections of preprints and papers at about 147,000 in 2025 alone; many
  credit real, already prominent scholars \parencite{zhao2026citations}.
  \fig[0.27\textheight]{ch16-resolve}
  \begin{itemize}
    \item A short shell script in Chapter 16 resolves each DOI through
      \code{doi.org}; the script that checked this book's bibliography also
      compares the title, first author and year.
    \item Even a perfect match does not show that the source says what the
      sentence citing it claims. Only reading the source beside the sentence
      shows that.
  \end{itemize}
  \kp{Never cite a paper you have not opened. An AI-assisted search finds
    candidates; it does not replace reading them.}
  \source{Figure from Chapter 16 of the book.}
\end{frame}

\begin{frame}{Reading, analysis and writing}
  \small
  \begin{itemize}
    \item With the PDF in front of it, a model can explain a method, translate,
      or answer questions about the paper; asked from memory, it invents
      details.
    \item It cannot judge whether the paper is any good. In 63 studies, models
      agreed with human screeners about nine times in ten, but their median
      accuracy in assessing risk of bias was 0.62, on only six assessments
      \parencite{laignelot2026systematicreviews}. A confident summary of a bad
      paper reads exactly like one of a good paper.
    \item For analysis, use models to write code, not to produce results. Fix
      the analysis plan before the agent starts and choose the statistical
      model yourself.
    \item Writing: help with structure, first drafts and proofreading; the
      authors remain responsible for the content.
    \item The deeper issue is voice. When a tool helps us develop our own
      draft, writing remains part of the thinking process. Delegating the
      draft can remove that opportunity to learn.
  \end{itemize}
\end{frame}

\begin{frame}{Review: what a journal can ask}
  \small
  \begin{itemize}
    \item Nature Portfolio: AI cannot be an author, because it cannot take
      accountability. Authors who use it to help write declare this in the
      Methods and check the text \parencite{naturemethods2026responsibly}.
      Records a model helped to produce, deposited in GBIF or OBIS, should say
      so too.
    \item Reviewers may not upload a manuscript under review into a generative
      AI tool: unfair to the authors, and a breach of confidentiality. No
      exception for tools that run on your own machine.
    \item An estimated 6.5 to 16.9 per cent of the text in peer reviews for
      four AI conferences in 2023 and 2024 was substantially modified by a
      model \parencite{liang2024peerreview}.
    \item Automated checks (references, statistics, duplicated figures) need
      no generative model. A model's review is text shaped like a review,
      without the judgement behind it.
    \item IPBES asks for traceability: a generative summary that cannot show
      where each statement came from fails \parencite{krug2026ipbes}. Do not
      use AI-text detectors to judge people; make disclosure easy and normal.
  \end{itemize}
\end{frame}

% ---------------------------------------------------------------------------
\section{How to keep up (Chapter 17)}

\begin{frame}{What to learn, and how long it lasts}
  \scriptsize
  \begin{tabularx}{\textwidth}{XlX}
    \toprule
    \textbf{Layer} & \textbf{Stays true for} & \textbf{Worth learning?} \\
    \midrule
    Model names and versions & Weeks & No. Look them up. \\
    Benchmark scores & Months & Only as a first filter for which models to test on your own task. \\
    Prices and usage limits & Months & Only when buying. \\
    Tool names and options & 6--12 months & Lightly; read the documentation. \\
    Reference databases and taxonomic backbones & Months to years & Not to learn, but follow their release notes: a new release can change your results. \\
    Patterns (harness, RAG, agent loop) & Years & Yes. \\
    Failures (hallucination, sycophancy, injection) & Years & Yes, thoroughly. \\
    Data standards, provenance, reproducibility & Decades & Yes. This is what lasts. \\
    Ecological expertise & A career & Yes. It is what makes you useful here. \\
    \bottomrule
  \end{tabularx}
  \kp{Spend your learning time at the bottom of that table. Everything at the
    top is a lookup.}
  {\scriptsize Measured: of 50 widely cited findings re-tested on 20 models, 44
    per cent had expired by September 2026 and 34 per cent had weakened; what
    lasted was judgement, not capability \parencite{queen2026shelflife}.\par}
  \source{Table from Chapter 17 of the book.}
\end{frame}

\begin{frame}{People worth following, and where to look}
  \footnotesize
  \begin{columns}[T]
    \begin{column}{0.50\textwidth}
      \textbf{Careful rather than fast.} Everyone with three or more works
      in the bibliography, plus the author's picks; a selection:
      \begin{itemize}
        \item Melanie Mitchell (Santa Fe Institute): a careful sceptic of
          claims about AI capability.
        \item Yoshua Bengio (Mila and LawZero): read him alongside Mitchell,
          who reads the same events differently.
        \item Becker, Rush and Rein (METR): measuring what AI does to real
          work.
        \item Anil Seth (Sussex), Blaise Ag\"uera y Arcas (Google): the
          consciousness exchange.
        \item Donat Agosti (Plazi): biodiversity knowledge out of the
          published literature.
        \item Nakagawa, Lagisz and Parker: how far ecological evidence can be
          trusted.
      \end{itemize}
    \end{column}
    \begin{column}{0.48\textwidth}
      \textbf{Sources.}
      \begin{itemize}
        \item arXiv and bioRxiv: search them when you have a question; do
          not try to follow them. Once a year, the AI Index.
        \item Careful writers who test things by hand and report what fails:
          Simon Willison, Stephen Turner.
        \item Documentation: Hugging Face, Ollama's library, your harness's
          own.
        \item Your domain community is more useful than the AI one: TDWG,
          ELIXIR, the GBIF and OBIS channels.
        \item A survey in Greece: most environmental omics researchers had
          the hardware; the gaps were expertise, training and staff
          \parencite{theodosopoulou2022greek}.
      \end{itemize}
    \end{column}
  \end{columns}
\end{frame}

\begin{frame}{A routine, and a small project}
  \fig[0.24\textheight]{ch17-routine}
  {\scriptsize\color{dim}About forty hours a year; more starts competing with
    doing ecology.\par}
  \medskip
  \fig[0.24\textheight]{ch17-small-project}
  {\scriptsize\color{dim}In the author's view a small project, such as a RAG
    system over your own reprints, is the fastest way to learn. PDF extraction
    is the hard part, chunking can matter as much as the model, and the model
    is confidently wrong in ways you can now recognise.\par}
  \source{Figures from Chapter 17 of the book.}
\end{frame}

% ---------------------------------------------------------------------------
\section{Epilogue: the wider frame (Chapter 18)}

\begin{frame}{Evolution and the biosphere}
  \paperslide{muller2026evolvableai}{\footnotesize
    \begin{itemize}
      \item When do AI systems whose components and learning rules can
        themselves evolve by natural selection become possible? The authors
        argue the scenario is underrated \parencite{muller2026evolvableai}.
      \item Treating deployed AI as an evolving population is a move
        ecologists are well placed to judge.
      \item \textcite{korecki2025biosphericai}: AI ethics centres on humans,
        at most on sentient beings; the biosphere should be the design
        target instead.
      \item Both connect to Vernadsky's transition from biosphere to
        \term{noosphere}, a sphere of thought.
      \item \textcite{vonneumann1955survive} asked the version that has aged
        best: has the world become too small, and its political units too
        unstable, for the technology being built in it?
    \end{itemize}}
  \source{\fullcite{muller2026evolvableai}}
\end{frame}

\begin{frame}{An inherited archive}
  \begin{itemize}
    \item David Krakauer's essay on the fossil logic of AI (Appendix C)
      compares extracting useful work from an inherited archive with
      extracting energy from fossil fuels.
    \item The difference matters: reading does not consume knowledge, but
      returning unchecked generated claims to the archive can damage its
      reliability. An interpretive analogy, not a measured ecological result.
    \item A measured counterpart: models trained repeatedly on generated text
      lose the rare parts of the original distribution first
      \parencite{shumailov2024collapse}; keeping enough real data in the mix
      can avoid this \parencite{seddik2024collapse}.
    \item For biodiversity, the author sees the same risk in generated
      occurrence records, trait values and phenotype assignments: once
      deposited in public databases, they become other people's evidence and
      training data.
    \item The practical implication: continued investment in collections,
      field observations and careful curation, which renew the evidence on
      which later models and researchers depend.
  \end{itemize}
\end{frame}

\begin{frame}{Consciousness: the shape of the disagreement}
  \begin{itemize}
    \item Like evolvable AI and biospheric ethics, consciousness turns on how
      built systems relate to living ones.
    \item \textcite{chandaria2026consciousness} organise competing theories
      around five things to examine: behaviour, computation, physical causal
      organisation, bodily self-maintenance, and interaction between an
      organism and its environment.
    \item Their numerical examples show how conclusions change with explicit
      assumptions; they are not measured probabilities that a particular model
      is conscious.
    \item For readers coming from biology, the framework makes the
      disagreement clear: do the functions suffice, or do the living processes
      that realise them matter too?
    \item The uncertainty runs both ways: attributing experience wrongly can
      divert concern from humans and other animals; overlooking it, if it
      exists, could conceal suffering.
  \end{itemize}
\end{frame}

\begin{frame}{The exchange in \emph{Behavioral and Brain Sciences}}
  \paperslide{seth2026naturalism}{\footnotesize
    \begin{itemize}
      \item Seth: consciousness may depend on the processes by which living
        systems maintain themselves. Unlikely in conventional digital AI; open
        for systems that reproduce relevant properties of brains or life
        \parencite{seth2026naturalism}.
      \item Richards and Ag\"uera y Arcas: living processes matter, but
        prediction, bodily regulation and self-production are themselves
        computational, so artificial consciousness is likely. A hypothesis,
        not an observation \parencite{richards2026computational}.
      \item Neither makes fluent conversation a decisive test.
      \item The author's judgement: keep the questions visible; assess
        research systems by their observable effects on data, people and the
        biosphere.
    \end{itemize}}
  \source{\fullcite{seth2026naturalism}}
\end{frame}

% ---------------------------------------------------------------------------
\section{The appendices}

\begin{frame}{What each appendix is for}
  \small
  \begin{description}
    \item[A. Glossary] The terms the book uses, from language models and
      agents, from biodiversity data and from the statistics of testing a
      model. Protein and genomic language models are outside its scope.
    \item[B. Checklists] Before starting an AI-assisted project; each agent
      session; testing a model or assistant; before submitting a manuscript;
      before reviewing one; the minimum contents of an institutional AI
      policy.
    \item[C. Tutorials] Online books, videos, newsletters, blogs, awesome
      lists and tool repositories, with names checked against the sources.
      Start with \emph{The Turing Way}, Stanford's CME 295 and EMBL-EBI's
      recorded webinar series on language models in bioinformatics; the
      Carpentries, TeSS, EMBL-EBI Training and the GBIF and OBIS courses age
      more slowly. If a link fails, prefix it with
      \code{https://web.archive.org/web/}.
    \item[D. State of play] Names, prices and limits as they stood between
      September and November 2026, each section carrying the date on which it
      was checked: machines and prices, models, harnesses and where each keeps
      its rules and memory, Ollama's defaults, the versions used in the
      workshop, subscription limits. The chapters
      keep the reasoning and the lasting names, and point here for the rest.
  \end{description}
\end{frame}

\begin{frame}{The checklists that matter most}
  \footnotesize
  \begin{columns}[T]
    \begin{column}{0.49\textwidth}
      \textbf{Before starting an AI-assisted project}
      \begin{itemize}
        \item[$\square$] Data classified: what may go to a hosted service,
          what may not. Sensitive localities kept off hosted services.
        \item[$\square$] Taxonomic backbone named, with its version or date;
          coordinate reference system recorded.
        \item[$\square$] \code{README.md} and \code{AGENTS.md} written;
          \code{data/raw/} read-only and checksummed; backup verified.
      \end{itemize}
      \textbf{Each agent session}
      \begin{itemize}
        \item[$\square$] Before: checkpoint committed; no credentials
          readable; task small enough to review in one sitting.
        \item[$\square$] After: whole diff read; new dependencies verified;
          run twice with identical output; prompt, model and version logged.
      \end{itemize}
    \end{column}
    \begin{column}{0.49\textwidth}
      \textbf{Before submitting a manuscript}
      \begin{itemize}
        \item[$\square$] Every DOI resolves and its title matches; every cited
          paper opened, and the claim matches what it says.
        \item[$\square$] AI use disclosed specifically: model, version, date,
          purpose, checks.
        \item[$\square$] Every field a model extracted has a validation sample
          with its error rates; machine-derived records flagged before
          deposit.
      \end{itemize}
      \textbf{Before reviewing a manuscript}
      \begin{itemize}
        \item[$\square$] Journal policy on AI in peer review checked.
        \item[$\square$] Manuscript not uploaded to any hosted service. Local
          model or none.
        \item[$\square$] The judgement of soundness is yours.
      \end{itemize}
    \end{column}
  \end{columns}
  \source{Appendix B of the book, which also has the testing and
    institutional-policy lists.}
\end{frame}

% ---------------------------------------------------------------------------
\begin{frame}{Takeaways}
  \begin{itemize}
    \item Treat everything an agent reads as untrusted, including your own
      data files if they came from elsewhere, and what tools say about
      themselves.
    \item What surrounds the model matters as much as the model: least
      privilege, a restricted directory, a sandbox with both limits, and no
      network when the task is local.
    \item Never cite a paper you have not opened. An AI-assisted search finds
      candidates; it does not replace reading them.
    \item The skills that last are the ones ecology already values: knowing
      where data came from, doubting a result that is too convenient, and
      recording what you did so that someone else can check it.
    \item Generated records deposited in public databases become other
      people's evidence and training data; collections, field observations and
      curation renew the evidence.
    \item It is legitimate to decide that an AI method does not suit a
      problem; \textcite{mishra2026notready} makes that case for ecology
      seriously and well. What is not reasonable is using these tools without
      checking, or staying out of the argument.
  \end{itemize}
\end{frame}

\begin{frame}{What to tell colleagues}
  \begin{quote}
    These systems are useful for transforming text and writing code, where
    you supply the content or can cheaply check the result. They are
    unreliable for facts: they invent citations and agree with you when you
    are wrong. The main change to daily work is agentic coding, which is
    closer to system administration than to programming. The main thing to
    get right is data organisation and provenance, which was worth getting
    right anyway. Run models locally where you can, verify everything that
    matters, and disclose what you used.
  \end{quote}
  \source{Chapter 17 of the book.}
\end{frame}

\begin{frame}[allowframebreaks]{References}
  \printbibliography[heading=none]
\end{frame}

\end{document}
